% ============================================================================
% EXEMPLO FICTÍCIO — Relatório Técnico de Extensão como Prática Profissional
% Curso Técnico Integrado em Informática (PPC 2012) — IFRN Campus Ceará-Mirim
% ============================================================================
%
% Este arquivo é o irmão de "relatorio-metalinguistico.tex": mesmo modelo,
% mesma estrutura de seções, mas com um projeto fictício preenchido do
% início ao fim, para servir de referência visual de como o relatório
% pronto deve ficar.
%
% ATENÇÃO: dados, números, nomes de pessoas e de escolas deste exemplo são
% inventados para fins didáticos. Não citar nem reutilizar como se fossem
% reais.
%
% Compila com pdflatex (rode duas vezes por causa do sumário/referências).
% ============================================================================

\documentclass[article,12pt,a4paper]{abntex2}

\usepackage[T1]{fontenc}
\usepackage[utf8]{inputenc}
\usepackage[brazil]{babel}
\usepackage{times}
\usepackage{indentfirst}
\usepackage{graphicx}
\usepackage{microtype}

\renewcommand{\maketitlehookb}{}%% sem título em língua estrangeira (sem abstract)
\titulo{Manutenção Preventiva e Corretiva em Laboratório de Informática de Escola Pública Municipal: um Relato de Extensão}
\autor{Rafael dos Santos Medeiros \and Ana Beatriz Lima Costa \and Prof.\ Camila Ferreira Duarte\footnote{Professora orientadora. Docente da área de Informática, IFRN Campus Ceará-Mirim.}}
\local{Ceará-Mirim -- RN}
\data{2026}
\instituicao{Instituto Federal de Educação, Ciência e Tecnologia do Rio Grande do Norte -- IFRN \\ Campus Ceará-Mirim}

\begin{document}
\maketitle

\begin{resumo}
Laboratórios de informática de escolas públicas municipais frequentemente
ficam parcialmente inoperantes por falta de manutenção técnica
especializada, o que reduz o uso pedagógico do espaço. Este relatório tem
como objetivo descrever a ação de extensão de manutenção preventiva e
corretiva realizada no laboratório de informática da Escola Municipal
Professora Iracema Dantas, em Ceará-Mirim/RN, executada por estudantes do
curso Técnico Integrado em Informática do IFRN. A ação consistiu em
diagnóstico individual de cada máquina, limpeza física, atualização de
sistema operacional, remoção de softwares maliciosos e configuração de
rede local, ao longo de seis visitas técnicas quinzenais. Do total de 18
computadores do laboratório, 13 estavam parcial ou totalmente inoperantes
no diagnóstico inicial; ao final da ação, 17 máquinas foram recolocadas em
pleno funcionamento, restando 1 computador com defeito de placa-mãe sem
reparo viável. A escola voltou a usar o laboratório em todas as turmas do
turno vespertino, o que não ocorria havia mais de um ano. Conclui-se que a
ação de extensão restabeleceu a infraestrutura de informática da escola
parceira e proporcionou aos estudantes do IFRN experiência prática direta
de manutenção de computadores em ambiente real de trabalho.

\vspace{\onelineskip}
\noindent
\textbf{Palavras-chave}: Manutenção de computadores. Extensão. Laboratório de informática. Escola pública.
\end{resumo}

\textual

\section{Introdução}

A Escola Municipal Professora Iracema Dantas, localizada na zona rural de
Ceará-Mirim/RN, possui um laboratório de informática com 18 computadores,
montado em 2019 com recursos de um programa federal de inclusão digital.
Sem quadro técnico próprio de TI, a escola dependia de visitas esporádicas
da Secretaria Municipal de Educação para qualquer manutenção, o que, em
2025, resultou em mais da metade das máquinas fora de uso e na suspensão
das aulas de informática básica previstas no calendário escolar.

O problema motivador desta ação de extensão foi justamente essa lacuna:
uma escola pública com equipamento já disponível, mas inoperante por falta
de manutenção técnica, e sem outro canal de suporte disponível no curto
prazo. Essa situação está diretamente relacionada aos conteúdos das
disciplinas de Organização e Manutenção de Computadores e Arquitetura de
Redes de Computadores, do núcleo tecnológico do curso, o que justifica a
ação como aplicação prática desses conhecimentos.

O objetivo geral deste projeto foi restabelecer o funcionamento do
laboratório de informática da escola parceira. Os objetivos específicos
foram: (a) diagnosticar individualmente cada uma das 18 máquinas; (b)
executar a manutenção preventiva e corretiva necessária em cada uma; e (c)
configurar a rede local do laboratório para acesso à internet.

Este relatório está organizado da seguinte forma: a seção 2 apresenta os
conceitos técnicos que fundamentam a ação; a seção 3 descreve a
metodologia e as atividades realizadas; a seção 4 apresenta os resultados
e o impacto na escola parceira; e a seção 5 traz as considerações finais.

\section{Fundamentação}

A manutenção preventiva de computadores consiste em procedimentos
regulares de limpeza física, verificação de temperatura dos componentes e
atualização de software, realizados antes que uma falha ocorra, com o
objetivo de prolongar a vida útil do equipamento. Já a manutenção
corretiva atua após a identificação de uma falha específica, como
travamentos frequentes, superaquecimento ou infecção por software
malicioso (TORRES, 2016).

Em laboratórios escolares compartilhados por muitos usuários, sem
supervisão técnica constante, é comum o acúmulo de programas indesejados e
a degradação do desempenho ao longo do tempo, o que reforça a necessidade
de manutenção periódica, e não apenas corretiva pontual (ALENCAR;
FERREIRA, 2019).

No âmbito do IFRN, a extensão é compreendida como via de interação entre a
instituição e a comunidade, capaz de proporcionar troca de conhecimento e
qualificação profissional a ambas as partes envolvidas. Os conceitos de
manutenção preventiva e corretiva, e de configuração de rede local,
tratados nesta seção, são conteúdos da disciplina de Organização e
Manutenção de Computadores e de Arquitetura de Redes de Computadores do
curso, o que situa esta ação como aplicação direta desses conteúdos em
contexto real.

\section{Metodologia e Descrição das Atividades}

Esta ação de extensão foi planejada em três etapas: diagnóstico,
execução da manutenção e configuração de rede, realizadas ao longo de seis
visitas técnicas quinzenais à Escola Municipal Professora Iracema Dantas,
no período de agosto a novembro de 2026, conforme o plano de atividades da
prática profissional deferido pela orientadora.

O público atendido foi a comunidade escolar da instituição parceira: 18
computadores do laboratório de informática, usados por aproximadamente 420
estudantes matriculados nos turnos matutino e vespertino, e por 22
professores da escola.

Na primeira visita, cada um dos 18 computadores foi identificado e
diagnosticado individualmente, registrando-se sintomas relatados pelos
professores (lentidão, travamento, não ligar) e o estado físico do
equipamento. Nas quatro visitas seguintes, foram executados: limpeza
física interna (remoção de poeira, verificação de coolers), atualização
do sistema operacional, remoção de softwares maliciosos identificados por
antivírus, e substituição de três memórias RAM danificadas, doadas pelo
IFRN a partir de equipamentos desativados. Na sexta e última visita, foi
configurada a rede local do laboratório, com reorganização do cabeamento
e verificação de acesso à internet em todas as máquinas recuperadas.

Foram usados: kits de ferramentas de manutenção do laboratório do IFRN,
software antivírus gratuito, mídia de instalação do sistema operacional
Linux Ubuntu 22.04 (escolhido pela escola por não exigir licenciamento), e
um roteador de reposição doado pelo IFRN para substituir o equipamento
anterior, que apresentava defeito intermitente.

\section{Resultados e Impacto}

A Tabela \ref{tab:resultados-ext} apresenta a situação do laboratório
antes e depois da ação de extensão.

\begin{table}[!ht]
  \centering
  \caption{Situação do laboratório antes e depois da ação}
  \label{tab:resultados-ext}
  \begin{tabular}{lcc}
    \hline
    Situação & Antes & Depois \\
    \hline
    Computadores em pleno funcionamento & 5 & 17 \\
    Computadores com defeito reparável  & 9 & 0  \\
    Computadores com defeito irreparável & 4 & 1  \\
    Total de computadores & 18 & 18 \\
    \hline
  \end{tabular}
  \\[0.3em]
  {\footnotesize Fonte: elaborado pelos autores (2026), a partir dos diagnósticos registrados nas visitas técnicas.}
\end{table}

Das 13 máquinas inoperantes ou parcialmente inoperantes no diagnóstico
inicial, 12 foram recuperadas — das quais 3 exigiram troca de memória RAM
e as demais foram resolvidas por limpeza, atualização de software ou
remoção de malware. Apenas um computador permaneceu com defeito
irreparável (placa-mãe queimada), cujo custo de reparo excedia o valor do
equipamento.

Segundo relato da direção da escola, colhido na última visita, o
laboratório voltou a ser usado em todas as turmas do turno vespertino, o
que não ocorria havia mais de um ano, e a coordenação pedagógica já
reincluiu aulas de informática básica no planejamento do próximo bimestre.
Esse resultado responde diretamente ao objetivo geral do projeto, que era
restabelecer o funcionamento do laboratório, e ao terceiro objetivo
específico, de configurar a rede local — as 17 máquinas recuperadas
tiveram acesso à internet confirmado ao final da última visita.

\section{Considerações Finais}

Este relatório descreveu a ação de extensão de manutenção do laboratório
de informática da Escola Municipal Professora Iracema Dantas, que
recuperou 17 dos 18 computadores do laboratório e restabeleceu o uso
pedagógico do espaço pela escola, respondendo ao objetivo proposto.

Do ponto de vista da formação técnica, a ação permitiu aos estudantes do
IFRN aplicar, em condições reais de trabalho e fora do ambiente controlado
do laboratório do próprio câmpus, os conhecimentos de manutenção de
computadores e configuração de rede estudados no curso, em diálogo direto
com uma necessidade concreta da comunidade — o que caracteriza a unidade
entre teoria e prática exigida pela prática profissional.

Como limitação, a ação não incluiu treinamento dos professores da escola
para pequenos reparos e manutenção básica, o que teria ampliado a
autonomia da escola entre visitas futuras. Como continuidade, recomenda-se
que o projeto seja renovado no próximo edital de extensão, incluindo uma
oficina de manutenção básica para um professor de referência da escola,
de modo a reduzir a dependência de novas visitas técnicas externas.

\postextual

\section*{Referências}

\newcommand{\referencia}[1]{\par\noindent\hangindent=1cm #1\par\medskip}

\referencia{ALENCAR, R. M.; FERREIRA, T. S. Manutenção preventiva em
laboratórios de informática escolares: um estudo de caso em três escolas
públicas. \textit{Revista Técnica de Informática Aplicada}, Natal, v. 4,
n. 2, p. 45-58, 2019.}

\referencia{TORRES, Gabriel. \textit{Hardware}: manual completo. 7. ed. Rio
de Janeiro: Nova Terra, 2016.}


\end{document}
